\chapter{Deep Brain Stimulation}

\section*{Overview}
Deep brain stimulation (DBS) uses implanted electrodes and a programmable pulse generator to deliver electrical signals to selected brain regions.

\section*{Common Instruments and Equipment}
The exact setup varies by technique and setting. Common items may include:
\begin{itemize}[nosep,leftmargin=*]
  \item stereotactic frame or navigation system
  \item intracranial electrode and test-stimulation equipment
  \item implantable pulse generator, extension, and tunneling instruments
\end{itemize}

\section*{Why It Is Done}
It is used for selected patients with medication-refractory Parkinson disease, essential tremor, dystonia, and some other movement disorders.

\section*{How It Is Performed}
Using brain imaging and stereotactic guidance, electrodes are placed in a targeted brain structure and connected by extension wires to a generator implanted under the skin of the chest. Programming is adjusted after surgery.

\section*{Contraindications}
No single contraindication applies to every patient; the treating team assesses whether the procedure is suitable for the individual and the reason for treatment.

\section*{Precautions}
The care team reviews relevant health conditions, medicines, allergies, and procedure-specific risks, and explains any changes needed for the individual.

\section*{Risks and Follow-up}
Bleeding, stroke, infection, seizures, hardware problems, mood or speech changes, and stimulation-related effects are possible. Long-term programming, battery management, medication review, and specialist follow-up are required.

\section*{References}
\begin{enumerate}
  \item MedlinePlus. Deep Brain Stimulation. U.S. National Library of Medicine; 2025.
  \item Current Medical Diagnosis and Treatment. McGraw-Hill; 2025.
\end{enumerate}
\clearpage
